\documentclass[10pt,a4paper]{amsart}
\usepackage[margin=19mm]{geometry}
\usepackage{amsmath,amssymb,mathtools}
\usepackage[scr=rsfs,cal=euler]{mathalfa}
\usepackage{xfrac}
\usepackage[unicode,hidelinks,bookmarks=false]{hyperref}
\newcommand{\ie}{i.e.\ }
\newcommand{\cf}{cf.\ }
\newcommand{\resp}{respectively\ }
\newcommand{\Subsection}{\subsection}
\providecommand{\nobreakdash}{\nobreak\hbox{-}}
\setlength{\emergencystretch}{2em}
\multlinegap0pt
\usepackage{mathtools}
\usepackage{amssymb}
\usepackage{xcolor}
\usepackage[shortlabels]{enumitem}
\usepackage{float}
\newtheorem{definition}{Definition}
\newtheorem{lemma}{Lemma}
\title{\textbf{Zero-Freeness of the Hard-Core Model with Bounded Connective Constant}}
\author{Yuan Chen, Shuai Shao, Ke Shi}
\date{Source: arXiv:2604.02746v1 (CC BY 4.0)}
\begin{document}
\maketitle

\noindent\emph{Source and licence.} \textbf{Zero-Freeness of the Hard-Core Model with Bounded Connective Constant}. Version arXiv:2604.02746v1 (CC BY 4.0), distributed by arXiv under the Creative Commons Attribution 4.0 International licence, http://creativecommons.org/licenses/by/4.0/, as recorded in the arXiv licence metadata for this exact version and shown on the abstract page https://arxiv.org/abs/2604.02746v1. Full licence notice: \texttt{licenses/arxiv-2604.02746-CC-BY-4.0.txt}.\par\smallskip

\noindent\textit{Editorial scope.} Counted unit: the article's lemma lem:P\_implies\_vanhove (source0.tex lines 882--893). The article's complete original proof of it is delivered with it (source0.tex lines 895--980, one \texttt{proof} environment of 86 lines). No mathematics is written, repaired or re-derived: the statement and the proof are the article's own lines, in the article's own notation, and no retained line is substituted or re-flowed.  Retained uncounted as the source-local context the proof needs: lines 833--874.  Nothing was omitted from any retained span, so the package carries no omission record.  No retained line contains a citation, so the wrapper carries no bibliography and no bibliography entry.  No retained line contains a figure environment, an image-inclusion command or drawing code, so no image is delivered and the package declares no asset dependency.  Editorial formatting only, in addition: an AMS \texttt{amsart} preamble that restates the article's own declarations and macros used by the retained lines, namely the environments \texttt{definition}, \texttt{lemma} on separate counters, together with the standard packages those lines need.  The retained environments print in this document as Definition 1, Lemma 1 in order of appearance, whereas the article prints the same items with the numbering of its own full document.  The complete native source of the article as distributed in the arXiv e-print for this exact version is bundled unchanged as \texttt{sources/197756/source0.tex}.
\begin{definition}[Assumption $\mathcal{A}$]\label{def:assumptionA}
Let $G=(V,E)$ be a connected, locally finite, infinite graph.
We say that $G$ satisfies \emph{Assumption $\mathcal{A}$} if there exist an integer $d\ge 1$
and a subgroup $\Gamma\le \mathrm{Aut}(G)$ such that the following hold:

\begin{enumerate}
% \item \textbf{Translation group.}
% There is a group isomorphism $\Gamma \cong \mathbb{Z}^d$.
% (Equivalently, $\Gamma$ is generated by $d$ commuting automorphisms of infinite order.)

\item \textbf{Translation group.}
There exist an integer $d\ge 1$ and automorphisms $\tau_1,\dots,\tau_d\in \mathrm{Aut}(G)$ such that
(i) $\tau_i\tau_j=\tau_j\tau_i$ for all $i,j$, and
(ii) the map $\mathbb{Z}^d\to\Gamma$ given by
\[
(n_1,\dots,n_d)\longmapsto \tau_1^{n_1}\cdots \tau_d^{n_d}
\]
is a group isomorphism. Equivalently, $\Gamma=\langle\tau_1,\dots,\tau_d\rangle=\{\tau_1^{n_1}\cdots\tau_d^{n_d}:(n_1,\dots,n_d)\in\mathbb{Z}^d\}$ is a free abelian group of rank $d$
(i.e., isomorphic to $\mathbb{Z}^d$).


\item \textbf{Free.}
The action of $\Gamma$ on $V$ is free, i.e.
\[
\forall v\in V,\quad \bigl(\gamma v = v \ \Rightarrow\  \gamma = e\bigr),
\]
where $e$ denotes the identity element of $\Gamma$.

\item \textbf{Cocompact.}
The action of $\Gamma$ on $V$ has finitely many vertex-orbits, i.e. there exists a finite set
$F\subset V$ such that
\[
V = \Gamma F := \{\gamma f:\ \gamma\in\Gamma,\ f\in F\}.
\]
Any such finite set $F$ is called a \emph{(finite) fundamental domain}.

% \item \textbf{$\Gamma$-invariant model parameters (optional, model-dependent).}
% Whenever a statistical-mechanics model is considered on $G$, we assume its parameters are
% $\Gamma$-invariant; for example, for the hard-core model this holds when the fugacity is constant,
% and more generally one may require $w_{\gamma v}=w_v$ for all $v\in V$ and $\gamma\in\Gamma$.
\end{enumerate}
\end{definition}

\begin{lemma}[Assumption $\mathcal{A}$ implies existence of a F\o lner-van Hove sequence]\label{lem:P_implies_vanhove}
Let $G=(V,E)$ be a connected, locally finite, infinite graph. Assume that $G$ satisfies
Assumption $\mathcal{A}$ from Definition~\ref{def:assumptionA}: namely, there exist an integer $d\ge 1$
and a subgroup $\Gamma\le \mathrm{Aut}(G)$ such that $\Gamma\cong \mathbb{Z}^d$, the action of $\Gamma$
on $V$ is free, and there exists a finite fundamental domain $F\subset V$ with $V=\Gamma F$.
Then $G$ admits a F\o lner-van Hove sequence $(\Lambda_n)_{n\ge 1}$ in the following:
for every fixed $r\in\mathbb{N}$,
\[
\frac{|\partial^{(r)}\Lambda_n|}{|\Lambda_n|}\longrightarrow 0 \qquad (n\to\infty),
\]
where $\partial^{(r)}\Lambda:=\{v\in \Lambda:\ \mathrm{dist}_G(v,V\setminus \Lambda)\le r\}$.
\end{lemma}

\begin{proof}
\textbf{Step 1: Reduce to a fundamental domain consisting of orbit representatives.}
Since $V=\Gamma F$ and $F$ is finite, $V$ is a union of finite number of $\Gamma$-orbits.
Choose a subset $F_0\subseteq F$ containing exactly one representative from each $\Gamma$-orbit.
Then $F_0$ is finite and
\[
V=\bigcup_{f\in F_0}\Gamma f.
\]
Replacing $F$ by $F_0$, we may assume without loss of generality that $F$ consists of orbit
representatives.

\smallskip
\textbf{Step 2: Build a cell-aligned sequence of finite sets.}
Fix an isomorphism $\Gamma\cong \mathbb{Z}^d$ and identify elements $\gamma\in\Gamma$ with vectors
$\mathbf{n}=(n_1,\dots,n_d)\in\mathbb{Z}^d$. For $n\ge 1$ define the box
\[
B_n:=\{-n,-n+1,\dots,n\}^d\subset \Gamma
\]
and the corresponding finite vertex set
\[
\Lambda_n \;:=\; \bigcup_{\gamma\in B_n} \gamma F \;\subset\; V.
\]
Since the sets $\gamma F$ are pairwise disjoint (because $F$ contains orbit representatives and
the action is free), we have
\[
|\Lambda_n| \;=\; |F|\,|B_n| \;=\; |F|\,(2n+1)^d \xrightarrow[n\to\infty]{}\infty.
\]

\smallskip
\textbf{Step 3: Finite range of ``cell displacements'' across edges.}
Define the finite set of displacements
\[
D\;:=\;\Bigl\{\delta\in\Gamma:\ \exists f,f'\in F \text{ with } \{f,\delta f'\}\in E\Bigr\}.
\]
The set $D$ is finite because $F$ is finite and $G$ is locally finite (only finitely many edges
emanate from vertices in $F$). Let
\[
M \;:=\; \max_{\delta\in D}\|\delta\|_\infty \;<\;\infty,
\]
where $\|\cdot\|_\infty$ denotes the $\ell_\infty$-norm under the identification $\Gamma\cong\mathbb{Z}^d$.

\smallskip
\textbf{Step 4: Points deep inside the box are far from the complement.}
Fix $r\in\mathbb{N}$ and set $s:=rM$. If $\gamma\in B_n$ satisfies
$\mathrm{dist}_\infty(\gamma,\Gamma\setminus B_n)>s$, then for any $f\in F$ we claim that
\[
\mathrm{dist}_G(\gamma f,\; V\setminus \Lambda_n) \;>\; r.
\]
Indeed, moving along one graph edge changes the cell-coordinate by an element of $D$, hence by
at most $M$ in $\ell_\infty$-norm. Therefore any path of length $\le r$ starting at $\gamma f$
remains within cell-coordinates contained in $B_n$, so it stays inside $\Lambda_n$.

Consequently, the $r$-thick boundary is contained in the union of cells whose coordinates lie
within $\ell_\infty$-distance $\le s$ of the box boundary:
\[
\partial^{(r)}\Lambda_n \;\subseteq\; \bigcup_{\gamma\in \mathrm{bd}_n^{(s)}} \gamma F,
\qquad
\mathrm{bd}_n^{(s)}:=\{\gamma\in B_n:\ \mathrm{dist}_\infty(\gamma,\Gamma\setminus B_n)\le s\}.
\]
Hence
\[
|\partial^{(r)}\Lambda_n|
\;\le\;
|F|\,|\mathrm{bd}_n^{(s)}|.
\]

\smallskip
\textbf{Step 5: Boundary-to-volume ratio tends to zero.}
Since $\mathrm{bd}_n^{(s)} = B_n\setminus B_{n-s}$ for $n>s$, we have
\[
|\mathrm{bd}_n^{(s)}|
\;=\;
|B_n|-|B_{n-s}|
\;=\;
(2n+1)^d - (2(n-s)+1)^d
\;=\; O(n^{d-1}).
\]
Therefore,
\[
\frac{|\partial^{(r)}\Lambda_n|}{|\Lambda_n|}
\;\le\;
\frac{|F|\,O(n^{d-1})}{|F|(2n+1)^d}
\;=\; O\!\left(\frac{1}{n}\right)\xrightarrow[n\to\infty]{}0.
\]
Since $r$ was arbitrary, $(\Lambda_n)$ is a Følner--van Hove sequence.
\end{proof}

\end{document}
